\chapter{Catalog of NOXNA Extensions}
\label{ch:noxna-catalog}

Volume~I's design-philosophy chapter introduced \texttt{NOXNA} as CNA's marker convention ---
and its matching compile-time purity check --- for any declaration inside the XNA namespace
tree that is not part of the real XNA~4.0 API. This appendix is not, and could not
responsibly be, an exhaustive line-by-line listing of every \texttt{NOXNA}-tagged member in a
codebase of several hundred header files; that list changes with every commit, and the
project's own \texttt{NOXNA} build gate (Volume~I, Chapter~3) already checks it mechanically,
which a static list in a book cannot. What this appendix collects instead is every
\texttt{NOXNA} (or the closely related \texttt{EXT}-suffixed) extension named specifically,
by class and member, across both volumes of this book --- organized by area, as a reference
for the extensions this book itself considers significant enough to have explained in the
main text, with a pointer back to where each one is covered in full.

\section{Core framework (Volume~I)}

\begin{description}
\item[\cnaclass{Game}] \texttt{RunApplication}, \texttt{getTargetFPSProperty()}/%
  \texttt{getTargetMsFrameTimeProperty()}, static \texttt{fpsToMillisecondsPerFrame()} --- Ch.6.
\item[\cnaclass{GameWindow}] \texttt{GetNativeSdlWindowEXT()}, \texttt{IsBorderlessEXTProperty},
  \texttt{MinimizeEXT()}/\texttt{RestoreEXT()} --- Ch.6.
\item[\cnaclass{GraphicsDeviceManager}] the entire \texttt{PresentationMode} enum
  (\texttt{Letterbox}/\texttt{Overscan}/\texttt{Stretch}/\texttt{NativeBackBuffer}/%
  \texttt{FixedHeightDynamicWidth}) and its get/set accessors --- Ch.6.
\item[\cnaclass{Matrix}] \texttt{ToColumnMajor(float[16])}, the one graphics-motivated bridge
  method on an otherwise pure math type --- Ch.7.
\item[\cnaclass{Color}] two byte-based constructors and a commutative \texttt{float * Color}
  operator, both explicitly marked as conveniences absent from real XNA --- Ch.7.
\end{description}

\section{Graphics (Volume~I)}

\begin{description}
\item[\cnaclass{Effect}] \texttt{Apply()} itself is \texttt{NOXNA} --- real FNA has no public
  \texttt{Effect.Apply()}, only \texttt{EffectPass.Apply()} --- Ch.13.
\item[\cnaclass{SkinnedEffect}] \texttt{VertexColorEnabled}, added for glTF-imported meshes
  carrying a \texttt{COLOR\_0} attribute; no such property exists on real XNA's own
  \cnaclass{SkinnedEffect} --- Ch.13.
\item[\cnaclass{IGraphicsBackend}] \texttt{SetStringMarkerEXT()},
  \texttt{DebugSimulateContextLoss()}/\texttt{DebugRestoreContext()},
  \texttt{UpdatePresentationFormatEXT()} --- Ch.16.
\item[Custom shaders] The entire \texttt{ShaderEffect} runtime-HLSL-compilation path on the
  three Direct3D backends --- Ch.15, Ch.22.
\end{description}

\section{Input (Volume~II, Chapter~1)}

The largest concentration of \texttt{NOXNA}/\texttt{EXT} surface in the whole project. Whole
\texttt{NOXNA} subsystems: \cnaclass{CNA::Input::Clipboard}, \cnaclass{Joysticks},
\cnaclass{Sensors}, \cnaclass{Power}, \cnaclass{Haptics}, \cnaclass{InputDevices}. Per-class
\texttt{EXT} additions: \cnaclass{Keyboard}'s \texttt{GetKeyFromScancodeEXT}/%
\texttt{GetModStateEXT}/\texttt{GetScancodeNameEXT}; \cnaclass{Mouse}'s
\texttt{SetCursor}/\texttt{ClickedEXT}/\texttt{SetCaptureEXT}/\texttt{GetGlobalPositionEXT};
the entire \cnaclass{MouseCursor} class; \cnaclass{GamePad}'s \texttt{GetGUIDEXT}/%
\texttt{SetLightBarEXT}/\texttt{SetTriggerVibrationEXT}/\texttt{GetGyroEXT}/%
\texttt{GetAccelerometerEXT}/\texttt{GetPlayerIndexEXT}/\texttt{GetPowerInfoEXT}; and
\texttt{TextInputEXT}'s IME-candidate and input-type-hint additions.

\section{Devices and Sensors (Volume~II, Chapter~4)}

The entire \cnans{CNA::Devices} namespace is \texttt{NOXNA} in spirit --- \cnaclass{Camera},
\cnaclass{SystemTray}, \cnaclass{DisplayInfo}, \cnaclass{MessageBox}, \cnaclass{FileDialog},
\cnaclass{UrlLauncher}, \cnaclass{Locale}, \cnaclass{Clipboard}, \cnaclass{SystemInfo},
\cnaclass{PowerInfo} --- since none has any real XNA/WP7 precedent at all. Within the
faithful \cnans{Microsoft::Devices::Sensors} namespace: \cnaclass{SensorBase}'s
\texttt{TimeBetweenUpdatesChanged} event (a real, corrected mistagging finding --- it was
briefly and incorrectly marked as real XNA API before being fixed); \cnaclass{Compass}'s
\texttt{getStateProperty()} (added purely for symmetry with \cnaclass{Accelerometer}'s real
\texttt{State} property, which real XNA's \cnaclass{Compass} itself lacks);
\cnaclass{VibrateController}'s \texttt{Start(TimeSpan, float intensity)},
\texttt{getIsSupportedProperty()}, \texttt{getDeviceNameProperty()}, and
\texttt{StartLeftRight()}.

\section{Storage (Volume~II, Chapter~8)}

\cnaclass{StorageDevice}'s \texttt{SetAppNameEXT()}/\texttt{GetStorageRootEXT()} --- the two
methods the entire namespace's real, local-filesystem-backed behavior ultimately roots
itself in.

\section{Content and assets (Volume~I, Chapter~8)}

The \texttt{.cnj} content format itself is a \texttt{NOXNA} system in its entirety, alongside
the Avatar system's separately-kept \cnaclass{SkinnedModelEXT} content pipeline
(\texttt{.skinnedmodel.json}/\texttt{.skeleton.bin}/\texttt{.clip.bin}) and its own
\texttt{NOXNA} vertex/track/clip types (\cnaclass{VertexPositionNormalTextureSkinned},
\cnaclass{KeyframeEXT}, \cnaclass{BoneTrackEXT}, \cnaclass{AnimationClipEXT}) --- see Volume~II,
Chapter~7.

\section{Avatar (Volume~II, Chapter~7)}

The entire \cnaclass{SkinnedModelEXT} real-rendering extension --- including
\texttt{AvatarRenderer::EnableRealRenderingEXT()}, \texttt{DrawRealEXT()},
\texttt{SetAppearanceEXT()}, \texttt{AvatarAnimation::SetRealClipNameEXT}/%
\texttt{GetRealClipNameEXT}, and the CNA-invented \cnaclass{AvatarAppearanceEXT} tint struct
--- is \texttt{NOXNA} in its entirety, deliberately opt-in and structurally separate from the
faithful base Avatar API it sits alongside.

\section{GamerServices (Volume~II, Chapter~5)}

\cnaclass{Guide}'s two real overlay-UI render hooks,
\texttt{RenderPendingMessageBoxEXT()}/\texttt{RenderPendingKeyboardInputEXT()}, and their
matching headless test-simulation hooks.

\section{How to read this appendix}

Treat this list the same way this book has asked you to treat every other list of specifics
it has offered: as a map into the source, not a replacement for it. The project's own
\texttt{NOXNA} compile-time gate (build with the \texttt{NOXNA} CMake option enabled) is the
authoritative, mechanically-enforced check for whether a given declaration inside the XNA
namespace tree is correctly tagged --- this appendix exists to give a reader who has just
finished both volumes a single place to recall which specific extensions this book found
significant enough to explain, not to serve as a substitute for that build-time check against
the codebase's current, ever-changing state.
